% part: intuitionistic-logic
% chapter: introduction
% section: natural-deduction

\documentclass[../../../include/open-logic-section]{subfiles}

\begin{document}

\olfileid[hi]{int}{int}{ntd}

\olsection{प्राकृतिक निगमन}

$\FalseCl$ नियमों के बिना प्राकृतिक निगमन, अंतःप्रज्ञावादी तर्कशास्त्र की
एक मानक !!{derivation}-प्रणाली है। हम नियमों को यहाँ दोहराते हैं और BHK
व्याख्या के द्वारा उनकी प्रेरणा बताते हैं। प्रत्येक मामले में नियम हमें यह
निष्कर्ष निकालने देता है कि यदि पूर्वाधारों की रचनाएँ हैं, तो निष्कर्ष की भी
रचना है।

चूँकि प्राकृतिक निगमन की !!{derivation}यों में अविमोचित मान्यताएँ होती हैं,
इसलिए ऐसी किसी !!{derivation} को---मसलन \mbox{$\Gamma$} की
!!{undischarged} मान्यताओं से~$!A$ की व्युत्पत्ति को---ऐसा फलन मानना चाहिए जो
सभी $!B \in \Gamma$ की रचनाओं को~$!A$ की रचना में बदलता है। यदि किसी भी
!!{undischarged} मान्यता के बिना~$!A$ की कोई !!{derivation} है, तो BHK व्याख्या
के अर्थ में~$!A$ की रचना है। किंतु चर्चा के लिए, जहाँ आवश्यकता न हो वहाँ हम
~$\Gamma$ को नहीं लिखेंगे।

मान्यता~$!A$ स्वयं, !!{undischarged} मान्यता~$!A$ से~$!A$ की एक
!!{derivation} है। यह BHK व्याख्या से मेल खाता है: रचनाओं पर तत्समक फलन
~$!A$ की किसी भी रचना को~$!A$ की रचना में बदलता है।

\subsection{संयोजन}

\begin{defish}
\AxiomC{$!A$}
\AxiomC{$!B$}
\RightLabel{\Intro{\land}}
\BinaryInfC{$!A \land !B$}
\DisplayProof
\hfill
\begin{tabular}{l}
\AxiomC{$!A \land !B$}
\RightLabel{\Elim{\land}}
\UnaryInfC{$!A$}
\DisplayProof
\\[3ex]
\AxiomC{$!A \land !B$}
\RightLabel{\Elim{\land}}
\UnaryInfC{$!B$}
\DisplayProof
\end{tabular}
\end{defish}

मान लें कि हमारे पास क्रमशः $!A_1$ और $!A_2$ की रचनाएँ $N_1$, $N_2$ हैं।
तब हमारे पास $!A_1 \land !A_2$ की रचना भी है, अर्थात युग्म
$\tuple{N_1, N_2}$।

BHK व्याख्या में $!A_1 \land !A_1$ की रचना एक युग्म
$\tuple{N_1, N_2}$ है। अतः मान लें कि हमारे पास ऐसा युग्म है। तब हमारे
पास प्रत्येक संयोज्य की रचना भी है: $N_1$,~$!A_1$ की और $N_2$,~$!A_2$ की
रचना है।

\subsection{सशर्त}

\begin{defish}
  \AxiomC{$\Discharge{!A}{u}$}
  \DeduceC{$!B$}
  \DischargeRule{$\Intro{\lif}$}{u}
  \UnaryInfC{$!A \lif !B$}
  \DisplayProof
\hfill
  \AxiomC{$!A \lif !B$}
  \AxiomC{$!A$}
  \RightLabel{$\Elim{\lif}$}
  \BinaryInfC{$!B$}
  \DisplayProof
\end{defish}

यदि हमारे पास !!{undischarged} मान्यता~$!A$ से~$!B$ की कोई !!{derivation} है,
तो ऐसा फलन~$f$ है जो~$!A$ की रचनाओं को~$!B$ की रचनाओं में बदलता है। वही
फलन $!A \lif !B$ की रचना है। अतः यदि $\Intro\lif$ के पूर्वाधार की रचना,
~$!A$ की किसी रचना पर आश्रित है, तो निष्कर्ष $!A \lif !B$ की रचना है।

दूसरी ओर मान लें कि~$!A$ की रचना~$N$ और $!A \lif !B$ की रचना~$f$ है।
$!A \lif !B$ की रचना ऐसा फलन है जो~$!A$ की रचनाओं को~$!B$ की रचनाओं में
बदलता है। अतः $f(N)$,~$!B$ की रचना है; अर्थात $\Elim\lif$ के निष्कर्ष की
रचना है।

\subsection{वियोजन}

\begin{defish}
\begin{tabular}{l}
\AxiomC{$!A$}
\RightLabel{\Intro{\lor}}
\UnaryInfC{$!A \lor !B$}
\DisplayProof
\\[3ex]
\AxiomC{$!B$}
\RightLabel{\Intro{\lor}}
\UnaryInfC{$!A \lor !B$}
\DisplayProof
\end{tabular}
\hfill
\AxiomC{$!A \lor !B$}
\AxiomC{$\Discharge{!A}{n}$}
\DeduceC{$!C$}
\AxiomC{$\Discharge{!B}{n}$}
\DeduceC{$!C$}
\DischargeRule{\Elim{\lor}}{n}
\TrinaryInfC{$!C$}
\DisplayProof
\end{defish}

यदि हमारे पास~$!A_i$ की कोई रचना~$N_i$ है, तो उसे हम
~$!A_1 \lor !A_2$ की रचना~$\tuple{i, N_i}$ में बदल सकते हैं। दूसरी ओर,
मान लें कि हमारे पास~$!A_1 \lor !A_2$ की कोई रचना, अर्थात ऐसा युग्म
$\tuple{i, N_i}$ है जहाँ $N_i$,~$!A_i$ की रचना है; और ऐसे फलन $f_1$,
$f_2$ भी हैं जो क्रमशः $!A_1$, $!A_2$ की रचनाओं को~$!C$ की रचनाओं में
बदलते हैं। तब $f_i(N_i)$,~$\Elim\lor$ के निष्कर्ष~$!C$ की रचना है।

\subsection{असंगति}

\begin{defish}
  \AxiomC{$\lfalse$}
  \RightLabel{\FalseInt}
  \UnaryInfC{$!A$}
  \DisplayProof
\end{defish}

यदि हमारे पास !!{undischarged} मान्यताओं~$!B_1$, \dots, $!B_n$ से~$\lfalse$
की कोई !!{derivation} है, तो ऐसा फलन~$f(M_1, \dots, M_n)$ है जो $!B_1$,
\dots, $!B_n$ की रचनाओं को~$\lfalse$ की रचना में बदलता है। चूँकि
$\lfalse$ की कोई रचना नहीं है, इसलिए $!B_1$, \dots, $!B_n$ सभी की रचनाएँ
भी नहीं हो सकतीं। अतः $f$ में यह गुण भी है कि \emph{यदि} $M_1$, \dots,
$M_n$ क्रमशः $!B_1$, \dots, $!B_n$ की रचनाएँ हैं, \emph{तो}
$f(M_1, \dots, M_n)$,~$!A$ की रचना है।

\subsection{$\lnot$ के नियम}

चूँकि $\lnot !A$ को $!A \lif \lfalse$ से परिभाषित किया जाता है, इसलिए ठीक
कहें तो हमें~$\lnot$ के नियमों की आवश्यकता नहीं है। फिर भी यदि हम उन्हें
दें, तो वे इस प्रकार दिखेंगे:

\begin{defish}
\AxiomC{$\Discharge{!A}{n}$}
\noLine
\DeduceC{$\lfalse$}
\DischargeRule{\Intro{\lnot}}{n}
\UnaryInfC{$\lnot !A$}
\DisplayProof
\hfill
\AxiomC{$\lnot !A$}
\AxiomC{$!A$}
\RightLabel{\Elim{\lnot}}
\BinaryInfC{$\lfalse$}
\DisplayProof
\end{defish}


\subsection{\usetoken{P}{derivation} के उदाहरण}

\begin{enumerate}
\item $\Proves !A \lif (\lnot !A \lif \lfalse)$, i.e., $\Proves !A
  \lif ((!A \lif \lfalse) \lif \lfalse)$
  \begin{prooftree}
    \AxiomC{$\Discharge{!A}{2}$}
    \AxiomC{$\Discharge{!A \lif \lfalse}{1}$}
    \RightLabel{\Elim\lif}
    \BinaryInfC{$\lfalse$}
    \DischargeRule{\Intro\lif}{1}
    \UnaryInfC{$(!A \lif \lfalse)\lif \lfalse$}
    \DischargeRule{\Intro\lif}{2}
    \UnaryInfC{$!A \lif (!A \lif \lfalse) \lif \lfalse$}
  \end{prooftree}

\item $\Proves ((!A \land !B) \lif !C) \lif (!A \lif (!B \lif !C))$
  \begin{prooftree}
    \AxiomC{$\Discharge{(!A \land !B) \lif !C}{3}$}
    \AxiomC{$\Discharge{!A}{2}$}
    \AxiomC{$\Discharge{!B}{1}$}
    \RightLabel{\Intro\land}
    \BinaryInfC{$!A \land !B$}
    \RightLabel{\Elim\lif}
    \BinaryInfC{$!C$}
    \DischargeRule{\Intro\lif}{1}
    \UnaryInfC{$!B \lif !C$}
    \DischargeRule{\Intro\lif}{2}
    \UnaryInfC{$!A \lif (!B \lif !C)$}
    \DischargeRule{\Intro\lif}{3}
    \UnaryInfC{$((!A \land !B) \lif !C) \lif (!A \lif (!B \lif !C))$}
  \end{prooftree}

\item $\Proves \lnot(!A \land \lnot !A)$, i.e., $\Proves (!A \land (!A
  \lif \lfalse))\lif \lfalse$
  \begin{prooftree}
    \AxiomC{$\Discharge{!A \land (!A \lif \lfalse)}{1}$}
    \RightLabel{\Elim\land}
    \UnaryInfC{$!A \lif \lfalse$}
    \AxiomC{$\Discharge{!A \land (!A \lif \lfalse)}{1}$}
    \RightLabel{\Elim\land}
    \UnaryInfC{$!A$}
    \RightLabel{\Elim\lif}
    \BinaryInfC{$\lfalse$}
    \DischargeRule{\Intro\lif}{1}
    \UnaryInfC{$(!A \land (!A \lif \lfalse))\lif \lfalse$}
  \end{prooftree}

\item $\Proves \lnot\lnot(!A \lor \lnot !A)$, i.e., $\Proves ((!A \lor
  (!A \lif \lfalse))\lif \lfalse)\lif \lfalse$
  \begin{prooftree}\hskip-3em
    \AxiomC{$\Discharge{(!A \lor (!A \lif \lfalse))\lif \lfalse}{2}$}
    \AxiomC{$\Discharge{(!A \lor (!A \lif \lfalse))\lif \lfalse}{2}$}
    \AxiomC{$\Discharge{!A}{1}$}
    \RightLabel{\Intro\lor}
    \UnaryInfC{$!A \lor (!A \lif \lfalse)$}
    \RightLabel{\Elim\lif}
    \BinaryInfC{$\lfalse$}
    \DischargeRule{\Intro\lif}{1}
    \UnaryInfC{$!A \lif \lfalse$}
    \RightLabel{\Intro\lor}
    \UnaryInfC{$!A \lor (!A \lif \lfalse)$}
    \RightLabel{\Elim\lif}
    \insertBetweenHyps{\hskip -4em}
    \BinaryInfC{$\lfalse$}
    \DischargeRule{\Intro\lif}{2}
    \UnaryInfC{$((!A \lor (!A \lif \lfalse))\lif \lfalse)\lif \lfalse$}
  \end{prooftree}
\end{enumerate}

\begin{prop}
  यदि अंतःप्रज्ञावादी तर्कशास्त्र के प्राकृतिक निगमन में
  $\Gamma \Proves !A$, तो शास्त्रीय तर्कशास्त्र में भी
  $\Gamma \Proves !A$। विशेषतः, यदि $!A$ अंतःप्रज्ञावादी प्रमेय है, तो वह
  शास्त्रीय प्रमेय भी है।
\end{prop}

\begin{proof}
  प्रत्येक प्राकृतिक निगमन नियम शास्त्रीय प्राकृतिक निगमन का भी नियम है;
  अतः अंतःप्रज्ञावादी तर्कशास्त्र की प्रत्येक !!{derivation}, शास्त्रीय
  तर्कशास्त्र की भी !!{derivation} है।
\end{proof}

\begin{prob}
  निम्न !!{formula}ों की अंतःप्रज्ञावादी तर्कशास्त्र में !!{derivation}याँ
  दीजिए:
  \begin{enumerate}
    \item $(\lnot !A \lor !B) \lif (!A \lif !B)$
    \item $\lnot\lnot\lnot !A \lif \lnot !A$
    \item $\lnot\lnot (!A \land !B) \liff (\lnot\lnot !A\land \lnot \lnot !B)$
    \item $\lnot (!A \lor !B) \liff (\lnot !A \land \lnot !B)$
    \item $(\lnot !A \lor \lnot !B) \lif \lnot (!A \land !B)$
    \item $\lnot\lnot ( !A \land !B) \lif  (\lnot\lnot !A \lor
    \lnot\lnot !B)$
  \end{enumerate}
\end{prob}

\end{document}
